\scanpage{172}
% archive foliation: 171
et si $S$ est le spectre d'un anneau de Dedekind
\[
H^1_{C_2}\xrightarrow{\sim}H^1_{C_1}
\]
[d'où dans ce cas $H^1_{C_3}\xrightarrow{\sim}H^1_{C_2}\simeq H^1_{C_1}$].
Si $S$ est le spectre d'un anneau \deleted{local complet} \inserted{de valuation} discrète à corps résiduel alg. clos $V$, on aura de $=$
\[
H^1_{C_i}\unsure{S|\delta}\simeq H^1_{\deleted{\unsure{C_i}}}(\widehat V/V,-)
\qquad\text{pour }i=1,2,3
\]

Si $S$ est le spectre d'un anneau local \uline{complet} $V_{\unsure{\varphi}}$, alors
\[
\deleted{\unsure{H^1_{C_4}(\widehat V/V,G)\simeq}}
\]
\[
H^1_{C_5}\simeq H^1_{C_4}\simeq H^1_{C_3}\quad(\simeq H^1_{C_2})
\]
et on aura de plus
\[
H^1_{C_i}(V,-)\simeq H^1_{C_i}
(\underset{\substack{\downarrow\\\text{corps}\\\text{résiduel}}}{k},-)
\qquad(i=2,3,4,5)
\]
